% Revised manuscript, 2 October 2026. Build with make.
% Results are transcribed from the supplied updated draft; no experiments rerun.
\documentclass[pdflatex,sn-mathphys-num]{sn-jnl}
\usepackage[T1]{fontenc}
\usepackage{graphicx,amsmath,amssymb,booktabs,tabularx,longtable,array,xcolor}
\usepackage[section]{placeins}
\usepackage{needspace}
\usepackage{xr-hyper}
\externaldocument[S-]{supplementary}
\graphicspath{{figures/}}
\newcommand{\ORCID}[1]{\href{https://orcid.org/#1}{\textsuperscript{\textsc{ORCID}}}}
\definecolor{queryorange}{RGB}{145,80,0}
\newcommand{\authorquery}[1]{\par\smallskip\noindent\begin{minipage}{\linewidth}\color{queryorange}\textbf{[AUTHOR QUERY: #1]}\end{minipage}\par\smallskip}
\emergencystretch=1.5em
\raggedbottom
\setcounter{topnumber}{4}
\setcounter{bottomnumber}{3}
\setcounter{totalnumber}{6}
\renewcommand{\topfraction}{.9}
\renewcommand{\bottomfraction}{.85}
\renewcommand{\textfraction}{.05}
\renewcommand{\floatpagefraction}{.7}
\begin{document}
\title[VCF-RDFizer]{VCF-RDFizer: From a VCF File to Explicit, Verifiable, and Governable Semantic Genomic Data}
\author*[1,2]{\fnm{Elias} \sur{Crum}\ORCID{0009-0005-3991-754X}}\email{elias.crum@ugent.be}
\author[1]{\fnm{Bart} \sur{Buelens}\ORCID{0000-0001-7734-3747}}
\author[2]{\fnm{Ruben} \sur{Taelman}\ORCID{0000-0001-5118-256X}}
\author[1]{\fnm{Gökhan} \sur{Ertaylan}\ORCID{0000-0001-5602-6435}}

\affil[2]{\orgdiv{IDLab, Department of Electronics and Information Systems}, \orgname{Ghent University -- imec}, \orgaddress{\street{Technologiepark-Zwijnaarde 116}, \postcode{9052}, \country{Belgium}}}
\affil[1]{\orgdiv{Digital Biological Systems}, \orgname{Flemish Institute for Technological Research (VITO)}, \orgaddress{\street{Boeretang 200}, \postcode{2400}, \country{Belgium}}}



\abstract{
\textbf{Background:}
The Variant Call Format (VCF) is widely used, compact, and readily parsable by a variety of tools, but two key challenges remain for its wider use in clinical contexts. 
(1) Information crucial for interpretation is not represented explicitly within the format, which makes integrating VCF data with other data sources such as personal, clinical, or experimental data difficult. 
(2) Privacy and consent information is represented at file level and often not in a machine-readable way.
Here we demonstrate the representation of VCF in the Resource Description Framework (RDF) to address these challenges, and present VCF-RDFizer, a command-line tool that mediates the conversion of VCF into RDF that can be verified against its source, using the public VCF Core vocabulary.

\textbf{Results:}
We evaluate VCF-RDFizer using public genomic datasets, source-based validation, scaling experiments, and a linked workflow combining genomic observations with external knowledge and simulated access policies.
The results show that VCF data can be converted into an explicit RDF representation while retaining the information needed to check the transformed data against its source.
Complementary query- and constraint-based validation expose different classes of transformation errors, demonstrating the value of validating both data content and graph structure.
The resulting RDF representations also support queries that combine genomic observations with external variant knowledge, population information, and requester-specific policies while reproducing the results of an equivalent conventional VCF workflow.
These capabilities come with additional conversion and indexing costs: semantic processing is therefore most advantageous for workflows that repeatedly reuse, integrate, or govern genomic data, whereas native VCF tools remain better suited to many direct or narrowly targeted analyses.

\textbf{Conclusions:}
VCF-RDFizer supports verifiable conversion and reusable, policy-aware research queries; it is not a universal replacement for VCF. 
Its value depends on the relationships needed, the validation coverage achievable, and the cost of conversion and indexing.
}
\keywords{Genomic variation, Variant Call Format, RDF, SPARQL, Semantic validation, Data-use policies}
\maketitle

\section{Background}\label{sec1}
The Variant Call Format (VCF) is widely used to represent and exchange genomic variation data in research and clinical sequencing, including in pharmacogenomics and precision medicine use cases~\cite{vcf_spec_2024,pharmcat,souche_recommendations_2022}.
Variant data are increasingly interpreted alongside electronic health records, external biomedical resources, sample metadata, and governance constraints~\cite{gottesman2013emerge,allofus2024genomic}.
The integration of these heterogeneous data establishes links between variants and external knowledge, but VCF does not support representing these links in a queryable or machine-interpretable way.
Additionally, genomic data are highly sensitive personal data, which the EU General Data Protection Regulation (GDPR) treats as a special category~\cite{gdpr2016,shabani2018gdpr}; access controls and data governance are therefore necessary for practical use.
Such controls are currently applied per file, which for a VCF means a whole genome or cohort at once.
Yet genomic regions, and the medical annotations linked to them, differ in how sensitive they are, so reuse would benefit from access controlled per region or per variant, which VCF does not support.
Finally, questions over VCF files are usually answered programmatically, with command-line utilities such as bcftools \cite{BCFtools} or with custom scripts or pipelines, rather than declaratively over a shared, reproducible model.

These limitations become particularly apparent when genomic observations need to participate in workflows beyond variant calling. 
Clinical variant interpretation, for example, requires observations to be repeatedly reconciled with evolving population frequencies, clinical assertions, phenotypes and other evidence, and systematic reanalysis can change diagnoses as these external resources evolve~\cite{richards2015standards,hiatt2018systematic,deignan2019reevaluation}. 
Pharmacogenomic decision support similarly requires genomic results to be combined with medications, patient context and continually maintained clinical recommendations; however, genomic information remains poorly integrated with electronic health records and standards-based exchange remains uncommon in implemented genomic decision-support systems~\cite{wake2022pgx,johnson2024genetically}. 
Finally, secondary use of genomic data requires its interpretation to be coupled with the conditions under which that data may be used. 
Variation in consent language and data-use restrictions has motivated standards such as the GA4GH Data Use Ontology and Machine-Readable Consent Guidance precisely because these conditions are otherwise difficult to interpret and evaluate computationally~\cite{lawson2021duo,rehm2021ga4gh}.

None of these tasks is precluded by VCF, and mature application-specific pipelines already support them. 
Rather, each requires relationships that extend beyond the genomic file itself and are therefore commonly reconstructed through annotations, external databases, and application logic. 
An RDF representation provides an alternative in which genomic observations can remain distinguishable from their interpretations while being linked through globally identifiable resources to changing biomedical knowledge, clinical context, and machine-readable usage policies. 
This motivates representing VCF not as a replacement for its compact exchange format, but as a semantic layer for workflows in which genomic data must be integrated, queried, and governed together with external information.

Consider a question that combines a participant's variants, ClinVar~\cite{landrum2018clinvar} classifications, population allele frequencies, and the rules governing a requester's access. 
Conventional tools can answer it. 
However, the workflow must make assumptions about variant identity, annotation meaning, and permitted release explicit somewhere. 
Semantic-genomics initiatives already connect variant observations to curated resources and governance frameworks \cite{togovar_2022,sphn_semantic_2023}. 
The question here is whether a reusable representation that can be verified against its source can support that work without hiding the cost of obtaining it.

The Resource Description Framework (RDF) represents relationships as triples whose resources can be identified by Internationalized Resource Identifiers (IRIs) \cite{rdf,hogan_knowledge_graphs_2021}. 
Linked Data practice uses those identifiers to connect independently managed resources \cite{W3CLinkedDataWiki,bio2rdf_2008,openbiolink_2019,sib_2024}; SPARQL queries their relationships \cite{sparqllang}; and the Open Digital Rights Language (ODRL) expresses policies about their use \cite{odrl_2018}. 
Neither RDF nor a policy annotation enforces privacy by itself. 
A policy-aware workflow also needs an evaluator that decides each request, a defined release boundary, and deployment controls that enforce the decision~\cite{xacml3_2013}.

Existing VCF-to-RDF converters include TogoVar, JVarkit, SPARQLing Genomics, and BioInterchange \cite{togovar_2022,JVarkit_2015,sparqling_genomics_2020,sparqling_genomics_software,baran2015gfvo}. 
They differ in their target models and treatment of file metadata and sample data; the wider documentation-based comparison is in Supplementary Table~\ref{S-tab:intro-tool-comparison}. 
Native frameworks (non-RDF) such as Hail, GEMINI, and TileDB-VCF also integrate variant annotations \cite{hail_software,paila2013gemini,tiledb_vcf_software}, while Beacon, htsget, and FHIR Genomics support discovery, retrieval, and clinical exchange \cite{rambla2022beacon,kelleher2019htsget,fhir_genomics_ig}. 
These are alternatives for particular tasks, not systems that require replacement by RDF.

Biomedical semantic tools provide useful precedents for combining transformation with assessment: SPHN Schema Forge generates schemas and validation artifacts, FAIR-Checker combines queries and shapes to assess metadata, and a FHIR RDF enrichment pipeline evaluates the information needed by its intended workflow \cite{toure2025forge,gaignard2023fairchecker,ansari2025etl}. 
For VCF, the corresponding challenge is to preserve source-dependent information, test the resulting representation, and demonstrate an integration task at a measured computational cost.

We present VCF-RDFizer, a Python- and Docker-based tool that mediates a VCF-to-RDF conversion and validation workflow. 
Existing converters already show that VCF can be converted to RDF; this study instead evaluates three capabilities together: conversion that can be verified against the source file, links to external knowledge maintained separately from the converted data, and policy-aware release of the result. 
Three questions organize the evaluation: 
\textbf{RQ1}, which aspects of VCF information are preserved and verifiable; 
\textbf{RQ2}, what the explicit representation enables in a linked, policy-aware workflow; and 
\textbf{RQ3}, under which workloads and representation choices its costs are practical. 
The vocabulary was developed separately \cite{crum_vcfcore_2027}; the representation decisions needed to assess the tool are specified here.

\section{Implementation}\label{sec2}
\subsection{Representation contract}\label{subsec2.1}
VCF-RDFizer generates a knowledge graph that represents the source VCF document, not only an interpreted variant. 
The conversion is built using the VCF Core Vocabulary~\cite{crum_vcfcore_2027} whose RDF output separates file and header assertions, variant records, and sample/genotype assertions (Table~\ref{tab:contract}). 
The vocabulary aligns with FALDO for coordinates, the Sequence Ontology for variant types, HERO for genomic constructs, and features of the GA4GH Variation Representation Specification (VRS) \cite{faldo_2016,seq_ont_2005,hero-genomics_2025,ga4gh_varschema_2021,ga4gh2025VRS}. 
These alignments connect VCF-specific structure to existing models; they do not replace the need to preserve the source's interpretation context.

Deterministic IRI templates identify resources by source file and record position. 
The same biological alteration in two files therefore remains two observations, each with its own quality values, annotations, and provenance. 
A separate linkset can connect both to a shared variant identifier (Figure~\ref{fig:workflow}). 
File-scoped observation identity and cross-file variant identity serve different purposes and are not collapsed.

\begin{table}[!htbp]
\caption{
  Representation contract and assessment boundary. 
  The complete capability matrix is in Supplementary Table~\ref{S-tab:capabilities-summary}; Q1--Q13 are specified in Supplementary Table~\ref{S-tab:semantic-queries}, and shapes are the Shapes Constraint Language (SHACL) constraints of Section~\ref{subsec2.2}.}\label{tab:contract}
\centering\small
\setlength{\tabcolsep}{4pt}\renewcommand{\arraystretch}{1.15}
\begin{tabularx}{\textwidth}{@{}>{\raggedright\arraybackslash}p{0.20\textwidth}>{\raggedright\arraybackslash}X>{\raggedright\arraybackslash}X@{}}
\toprule
\textbf{Information} & \textbf{Representation} & \textbf{Assessment and limit}\\
\midrule
File and header & Versioned file, declarations, and metadata linked to their values & Q7--Q10 and shapes; header attribute-value agreement needs the deeper shapes.\\
\addlinespace
Record and alleles & Core fields, ordered alleles, typed values, and explicit missingness & Q1--Q4, Q11, and shapes; the digest is a summary, not exhaustive equality.\\
\addlinespace
INFO & Raw string plus optional typed, structured values & Q12 and structural checks; preservation differs from interpreting every annotation.\\
\addlinespace
Samples and FORMAT & Expanded sample-call resources or condensed, sample-ordered vectors & Q5--Q6 and Q13; individual condensed values require decoding.\\
\addlinespace
Version-dependent constructs & VCF 4.1--4.5 file classes; symbolic and newer fields carried through & Fixture coverage, not full interpretation of all constructs; local-to-global allele resolution is not implemented.\\
\addlinespace
Links and policies & Separate linksets and policy-defined release views & Evaluated independently of conversion; identity assumptions and simulated policies bound the use case.\\
\bottomrule
\end{tabularx}
\end{table}

Sample data can be represented in two different profiles. 
\emph{Expanded} makes sample calls and their FORMAT values directly addressable graph resources. 
\emph{Condensed} stores sample calls and FORMAT values in separate ordered vectors, retaining the values but requiring their positions to be decoded. 
Both retain file, header, and record context. 
For $V$ records, $S$ samples, and $F$ FORMAT keys, the repeated sample structure scales as $\mathcal{O}(S+VSF)$ triples expanded and $\mathcal{O}(S+VF)$ condensed, assuming bounded ploidy and value arity. 
Condensed representations retain the $VSF$ values inside literals: fewer triples do not mean that the sample information disappears. 
The detailed model and derivation are in Supplementary Section~\ref{S-sup:model}.

\begin{figure}[!htbp]
\centering
\includegraphics[width=\textwidth]{fig-workflow.pdf}
\caption{
  \textbf{Source observations remain separate while their links can be reused.} 
  Conversion creates file-scoped records with header and sample context. Separate linksets connect eligible alleles to a shared identifier under the stated reference and normalization assumptions. 
  The evaluated workflow joins these links to ClinVar and, for the cohort, a separate frequency file. 
  Policies select a requester's release view; source and policy checks assess that view before querying. 
  HDT and COTTAS archives (dashed branch) are optional: QLever builds its own index from the N-Triples, so neither archive is needed to query.
  }\label{fig:workflow}
\end{figure}

\subsection{Conversion, validation, and extension points}\label{subsec2.2}\label{subsec2.4}\label{subsec:validation}\label{subsec:plugins}
The Python command-line tool separates conversion, optional representation construction, validation, and packaging. 
Record, header, and file metadata become tab-separated logical sources for RDF Mapping Language (RML) mappings, executed by RMLStreamer \cite{rml_2024,rmlstreamer_2022}. 
The remaining triples (typed INFO and FORMAT values, QUAL, header attributes, and the sample layers) are written by direct emitters: Python functions in the tool that write N-Triples themselves, because an RML mapping cannot choose a literal's datatype row by row (a missing QUAL, for instance, must be typed as a null value rather than a decimal). 
This is a hybrid implementation: RML produces the nine-triple record backbone, 5.3\% of the HG005 graph, while emitters produce the remaining 94.7\%. 
Editing mappings alone therefore does not change the complete representation.

Parts are appended to plain or gzip-framed N-Triples and deleted after appending. 
Optional Header--Dictionary--Triples (HDT) and Parquet-based COTTAS outputs are built incrementally, without keeping a full decompressed duplicate \cite{hdt_2013,cottas_2025}; gzip and Brotli packages are offered as well for transport or archival tasks \cite{brotli_2019}. 
Before source RDF is removed, native readers decode each selected representation and compare its triple count with the source. 
This checks decoding and counts, not all values. 
QLever builds its own query index; its measurements must be distinguished from engines that query HDT or COTTAS directly \cite{qlever_2017,Taelman2018Comunica,comunica_cottas_2026}. 
Construction details are in Supplementary Section~\ref{S-sup:construction}.

Readable output is not necessarily faithful output, so the validation layer checks the graph against the file it came from. 
It asks the graph thirteen questions (Q1--Q13) in SPARQL and answers each question a second time, independently, from the source VCF with the cyvcf2 parser, using bcftools for exact FILTER text \cite{vcf_intro,cyvcf2_2017,BCFtools}. 
This parser-derived answer is the \emph{oracle}: both answers are normalized and must be exactly equal, and any missing, extra, or differing row is reported as a mismatch. 
Q1--Q6 summarize records, alleles, and genotypes (for example, records per chromosome window or calls per genotype class); Q7--Q10 compare file metadata and the inventories of predicates and classes the graph should contain; and Q11--Q13 compare the values themselves. 
Rather than comparing every value pairwise, Q11--Q13 hash each value with SHA-256 together with the IRI of the record it belongs to, and both routes count values per bucket of the hash's first byte (256 buckets). 
A changed value moves to another bucket and is detected unless another change exactly offsets it, so these digests are a strong summary, not a collision-free, value-by-value comparison. 
Fourteen preflight queries and SHACL constraints provide additional checks \cite{shacl2017}. 
The default shapes assess cardinality and datatypes; two deeper profiles assess uniqueness and value agreement. 
Evaluation execution limits and measured sensitivity are reported separately.

Beyond conversion and validation, two plug-ins work on the converted graph without modifying it: a linker that connects variants to shared identifiers, and a policy plug-in that releases requester-specific views. 
Linking leaves the converted graph unchanged: links are written to a separate linkset. 
The three linker tiers rewrite existing identifiers, join digest-pinned reference data, or confirm identifiers through a rate-limited live service. 
The evaluated SPDI (Sequence Position Deletion Insertion) linker minimally trims explicit alleles and harmonizes contig names to RefSeq accessions \cite{holmes2020spdi}; a second linker records Ensembl gene overlaps. 
This establishes equality joins under the tested normalization assumptions, not canonical identity across all equivalent variant expressions.

The policy plug-in evaluates ODRL rules targeting files or selections declared in SPARQL. 
Its profile implements default-deny and deny-overrides, with purposes drawn from a subset of the Data Use Ontology (DUO) \cite{odrl_2018,lawson2021duo}. 
An ownership rule states what withholding a resource also removes: a record's dependent calls and alleles, or a file's header. 
Each release view is checked for policy-digest agreement, permitted membership, absence of dangling references, and record equality with a VCF-text-derived oracle. 
This is policy-aware release, not anonymization or a complete security deployment; obligations are recorded but not enforced.

\subsection{Evaluation methods}\label{sec:design}\label{subsec2.3}
\subsubsection{RQ1: Which aspects of VCF information are preserved and verifiable?}
The public corpus spans five Personal Genome Project donations, an HG002 precisionFDA/Genome in a Bottle (GIAB) challenge file, the 2,504-sample 1000 Genomes chromosome-20 call set, HGSVC2 structural variants, and the HG004 and HG005 GIAB benchmarks \cite{personal_genome_proj,nist_gib_2025,thousand_genomes_2015,hgsvc_2021,giab_v421_2022}. 
The corpus experiment processes the first 250,000 records of eight inputs and HG005 whole; the large cohort is excluded from the expanded-profile corpus run and assessed through controlled slices. 
The inputs are listed in Supplementary Table~\ref{S-tab:test-corpus}.

Paired validation covers the difficult-input fixtures, a 10,000-line (9,986-record) single-sample fixture under both sample profiles and three RDF artifacts, and slices up to 100,000 HG005 records. 
A separate follow-up tests 250,000 records from the consumer genome \texttt{NG131FQA1I}. 
This run is reported separately because it ran after the campaign, with a later build, on QLever without shapes. 
A fault-injection experiment tests 113 deliberate corruptions with queries alone, with the default shapes, and with all shape profiles; unmutated controls test whether the validator rejects valid inputs. 
Separate retrieval runs compare the same thirteen answers at whole-genome scale without SHACL. 
The specification, exclusions, and oracle dependencies are in Supplementary Section~\ref{S-sup:validation}.

\subsubsection{RQ2: What does the explicit representation enable in a linked, policy-aware workflow?}\label{subsec:usecase-design}
The application question is: \emph{which participant--variant--gene matches connect an observed allele to a ClinVar classification in the fixed 81-gene American College of Medical Genetics and Genomics secondary-findings list (ACMG SF v3.2), and which may each requester receive?} 
The benchmark uses this frozen gene list, not a claim about the latest recommendations \cite{miller2023acmg}. 
It accepts ClinVar records with at least one review star and includes all admitted classifications. 
Pathogenic/likely pathogenic (P/LP) matches are counted separately; the workflow does not implement the additional criteria for reporting clinical secondary findings.

The RDF route links converted genomes to ClinVar and gene resources, materializes requester views, and queries each view with SPARQL. 
The baseline harmonizes contigs, annotates variants with bcftools, and evaluates the same case definition in a separate Python implementation. 
Sorted result tuples must agree by participant, gene, contig, position, and alleles before timings are compared. 
Counts are matches, not unique participants. 
Shared preprocessing and case definitions limit the independence of the comparison even though the downstream implementations are separate.

Three arms test integration, breadth, and scale: five genomes from three sources; 104 unrelated high-coverage 1000 Genomes participants, four per population \cite{byrska2022highcoverage}; and HG005 whole. 
The first two use normalized, split records within the gene spans. 
The cohort removes repeated panel-level INFO from each participant and represents the required frequencies once in a separate sites-only graph. 
The whole-genome answer is also checked against HG005's restricted-input answer. 
ClinVar is fixed at its 2026-09-13 release \cite{clinvar_vcf}.

All consents are simulated. 
Requesters represent clinical care, disease-specific cardiovascular research, and general research. 
File-level permissions, withdrawals, and a rule restricting records in 28 cancer-predisposition genes determine release. 
Each participant is represented by a single-sample file; this does not test selective masking within a multi-sample cohort file. 
Four applied rule changes compare where configuration or execution code must change, not user performance or maintainability. 
Supplementary Section~\ref{S-sup:usecase} gives the complete protocol.

\subsubsection{RQ3: Under which workloads and representation choices are its costs practical?}\label{subsec:execution}
Record ladders use 10,000, 100,000 and one million HG005 records, with three replicates, plus the whole genome. 
Sample ladders hold 10,000 cohort records and INFO fixed while varying one to 2,504 samples under both profiles. 
We measure stage timings, resident memory, transient workspace, and retained artifact sizes separately. 
Whole-corpus measurements are single runs, not replicated estimates.

Retrieval compares individual questions, a thirteen-question batch, and five regional questions against scanned and tabix-indexed VCF access. 
Regional windows span 1\,kb--10\,Mb, with twenty windows per size and three replicates; indexed readers discard overlapping records whose POS lies outside the window so that all routes answer the same question. 
Setup is separate from query time. The principal machines have eight cores and 33.6\,GB RAM, with no concurrent benchmark jobs. 
The 143-cell base campaign uses v3.1.0; the use case converts with v3.2.0 and links and governs with pre-release plug-in builds; later retrieval, review-validation, and plug-in runs are separate experiment families. 
Supplementary Section~\ref{S-sup:repro} records versions, hosts, exclusions, and unresolved archival details rather than treating them as one release.

\section{Results}\label{sec3}
\subsection{RQ1: The graph agreed with its source wherever compared, but fault coverage depends on the checks enabled}\label{subsec3.5}
Figure~\ref{fig:validation} sets each validation result against what a faithful conversion and a sensitive validator should produce. 
In the base campaign, 76 validations completed on fixtures and slices of up to 17.1M triples. 
All 984 applicable comparisons between a SPARQL answer and its oracle answer were exactly equal; four were not applicable, because the zero-record fixture has no genotypes for Q5 and Q6 in either of its two artifacts. 
All 1,144 cross-check invariants also held, for example that transitions plus transversions equal the biallelic single-nucleotide variants. 
Twenty-five validations ran the same queries on all four SPARQL engines, and every engine returned the same answers; the other 51 used one engine. 
All 152 HDT and COTTAS artifacts decoded to exactly the source graph's triple count.

\begin{figure}[!htbp]
\centering
\includegraphics[width=\textwidth]{fig-validation.pdf}
\caption{
  \textbf{Validation results against the expected outcome.} 
  Each bar is the outcome that a faithful conversion and a sensitive validator should produce, filled to the observed result. 
  (a) On valid inputs, every answer should equal the one derived from the VCF: the base campaign's paired comparisons (four not-applicable comparisons excluded), the triple count of every decoded HDT and COTTAS artifact, the separate whole-genome retrieval run, and the consumer-genome follow-up. 
  (b) Every one of the 113 injected faults should be detected.
  }\label{fig:validation}
\end{figure}

These results answer different questions and have different limits. 
A query pass establishes equality of the specified summaries, and a shape pass establishes the enabled constraints; neither proves complete semantic equivalence. 
Decoding checks establish triple counts, not every value. 
The whole-genome retrieval run compared answers without SHACL, and neither it nor the follow-up run is added to the base campaign's denominator.

In the consumer-genome follow-up (250,000 records of \texttt{NG131FQA1I}, 58.2M triples), the three differing answers came from the validator, not the conversion: the oracle did not yet count the phase sets the converter emits, and QLever prints decimal QUAL values in canonical form. 
A direct comparison found every record field equal to the VCF text; Supplementary Section~\ref{S-sup:validation} gives the diagnosis.

Of the 113 injected faults, the queries alone detected 96 (Figure~\ref{fig:validation}b). 
The other 17, in ten fault classes, changed relationships or values that no query summary inspects, such as which allele a genotype refers to, a sample's index, phasing, and header attribute values. 
The default shape profile detected none of these 17; adding the two deeper profiles detected all of them, and both unmutated control graphs still passed. 
On a graph of about 2,000 triples, the two deeper profiles took 33.7 and 73.2\,s, against 1.7\,s for the default profile. 
The default profile also ran in 58 base validations, on graphs of up to 0.96M triples, and found no violations.

During development, these checks found three real defects, all since fixed. 
The predicate and class censuses (Q9 and Q10) showed that a header-only file omitted the samples its header declares; the shapes reported genotypes that referred to allele resources a raw-INFO conversion never emitted; and the policy plug-in's own check found a withheld file's header surviving in release views. 
A record count alone would have shown none of them. 
Version fixtures (VCFv4.1--4.5, and an undeclared version) converted, and nine of eleven difficult inputs converted and validated; a truncated gzip was refused, while a malformed sites-only fixture left that paired-coverage test incomplete. 
Supplementary Sections~\ref{S-sup:inputs}, \ref{S-sup:validation}, and~\ref{S-sup:issues} retain these outcomes, repeatability checks, and unresolved oracle limitations.

\subsection{RQ2: Shared identifiers enabled reusable joins and verified, requester-specific views}\label{subsec:usecase}\label{subsec3.7}\label{subsec:policy}
The RDF and bcftools routes returned identical requester-specific match sets in every arm (Figure~\ref{fig:usecase-matches}). 
Each release view passed the policy and source checks. 
HG005's whole-genome answer matched its restricted-input answer exactly: 1,496 participant--variant--gene tuples, not 1,496 people. The unrestricted P/LP subset contained zero matches in the five-genome and whole-genome arms and eight in the cohort. 
These are classifications retrieved under the case definition, not clinically adjudicated secondary findings.

\begin{figure}[!htbp]
\centering
\includegraphics[width=\textwidth]{fig-usecase-matches.pdf}
\caption{
  \textbf{The RDF and bcftools routes returned identical match sets for every requester.} 
  Counts are participant--variant--gene tuples, not people. 
  Bars are the RDF route and dots the bcftools route; they agree in every cell. 
  CC denotes clinical care, DS disease-specific (cardiovascular) research, and GRU general research; ``All'' applies no policy. 
  Triple counts are the genomic observations before links. 
  (c) joins the cohort's matches to the separate panel-frequency graph at panel allele frequency (AF) $<0.01$.
  }\label{fig:usecase-matches}
\end{figure}

The cohort's rare-variant query\footnote{\url{https://github.com/ecrum19/vcf-rdfizer-testing/blob/main/benchmarks/use_case/acmg/cohort/rare.rq}} shows the benefit of keeping shared information separate. 
In the 1000 Genomes files, every participant's records repeat panel-level INFO fields that describe all 3,202 panel samples rather than the participant; they made up 93\% of each participant's triples, approximately 857M triples across the cohort. 
The cohort instead drops them and stores the frequencies the question needs once, in a 68,635-site graph of 6.5M triples. 
The SPDI links then joined genomic observations, ClinVar, and panel frequencies without embedding those frequencies in each genome, and the baseline's position-and-allele lookup returned the same rare-variant matches. 
The reduction comes from storing shared information once, not from RDF compression.

Shared identity also changed where one release rule was expressed. 
Restricting HFE C282Y to clinical care required a policy edit without modifying the RDF policy engine; the baseline required a small code and data edit. 
On the change fixture, the identity-targeted rule also selected HG002's record written on contig \texttt{6} rather than \texttt{chr6}. 
Three other changes were similarly small data edits in both routes. 
This demonstrates reusable configuration for the tested change, not a general advantage in usability or lines of code. 
Full change-impact and linking results are in Supplementary Section~\ref{S-sup:usecase}.

Most of the linked workflow's cost is paid once, before the first question (Table~\ref{tab:usecase-cost}). 
For the whole genome, converting it and ClinVar took 54\,min and linking 8\,min; writing, checking, and indexing the clinical requester's view took another 2.5\,h, 68\,min of it checking the view. 
The run peaked at 13.1\,GB of memory and 34\,GB of disk. 
After that, each query took 205--209\,s. 
Query times were similar in all three arms (172--209\,s), which suggests that the fixed ClinVar join dominated this workload, not that SPARQL performance is independent of scale. 
% TODO: re-do this experiment with different numbers of things here (and some granular defined policies?)
In the whole-genome arm, the clinical view released the complete genome and the other two views were empty, so this arm shows that release scales for that policy pattern, not that selective masking does.

\begin{table}[!htbp]
\caption{
  Stage costs of the linked workflow. 
  Conversion and linking are paid once per arm, as totals over its files. 
  The other stages are paid per requester, as ranges across the requesters: writing the release view, checking it against the policies and the VCF text (for a non-empty view, this includes building the view's own index), building the query index, and running the query, three times with a cold cache. 
  The whole-genome query index holds 701.5M triples, including ClinVar and the links.
  }\label{tab:usecase-cost}
\centering\footnotesize
\setlength{\tabcolsep}{3pt}\renewcommand{\arraystretch}{1.15}
\begin{tabularx}{\textwidth}{@{}>{\raggedright\arraybackslash}X r r r r r r@{}}
\toprule
 & \multicolumn{2}{c}{\textbf{Once per arm}} & \multicolumn{4}{c}{\textbf{Per requester}}\\
\cmidrule(lr){2-3}\cmidrule(l){4-7}
\textbf{Arm} & \textbf{Convert} & \textbf{Link} & \textbf{View} & \textbf{Check} & \textbf{Index} & \textbf{Query}\\
 & \textit{min} & \textit{min} & \textit{min} & \textit{min} & \textit{min} & \textit{s}\\
\midrule
Five genomes & 3.3 & 1.4 & 0.53--0.65 & 0.35--0.62 & 2.2 & 172--181\\
Cohort of 104 & 31 & 22 & 0.98--1.33 & 3--5 & 2.4--3.3 & 176--202\\
Whole HG005, clinical view & 54 & 8.0 & 38 & 68 & 45 & 205--209\\
Whole HG005, empty views & \multicolumn{2}{c}{\textit{as above}} & 28 & 1 & 2.2 & $<1$\\
\bottomrule
\end{tabularx}
\end{table}

\subsection{RQ3: Costs are practical for repeated, linked questions, not for one-pass analyses}\label{subsec3.3}\label{subsec3.4}\label{subsec3.6}\label{subsec:scale}\label{subsec:regional}
\subsubsection{Sample granularity determines the amount of graph}
On the record ladder, triple production was approximately proportional to records at 170.5--171.9 triples per record. 
Mapping-stage resident memory remained 1.0--1.7\,GB; total wall time and disk grew with the data. 
This supports bounded-memory mapping, not a claim that every downstream validator or query engine is memory-bounded.

The sample ladder shows a more consequential design choice (Figure~\ref{fig:samples}). 
At 10,000 records and 2,504 samples, expanded produced 627.4M triples and a 4.10\,GB HDT, against 59.2\,MB of HDT in condensed form. 
The triple-count ratio was $432\times$, but the HDT-byte ratio was $69\times$. 
Condensed triple count grew only 0.9\% across the ladder while its HDT grew $5.8\times$, because the vectors lengthened. 
Its lower cost is useful only where applications can accept decoding rather than direct sample-call addressability. 
Neither profile was used to convert the complete 2,504-sample call set.

\begin{figure}[!htbp]
\centering
\includegraphics[width=\textwidth]{fig-samples.pdf}
\caption{
  \textbf{Sample representation trades addressability for size.} 
  The same 10,000 records are emitted with one to 2,504 samples. 
  (a) Expanded resources grow with sample calls; condensed vectors keep triple counts nearly constant. 
  (b) Both representations grow in bytes. 
  (c) At 2,504 samples the expanded/condensed ratio depends on the measure. 
  The $\times390$ annotation in (a) instead compares expanded output at the two ends of the ladder. 
  Medians are shown where replicated.
  }\label{fig:samples}
\end{figure}

\subsubsection{Optional artifacts should not be counted as mandatory conversion}
Space-optimized storage reduced transient workspace by $9.19\times$ on the 100,000-record input at a 5.2\% time cost, without changing the triples. 
In the corpus, COTTAS occupied 0.37--0.55 times the gzip-framed N-Triples and 0.27--0.37 times HDT, but took 1.5--3.0 times longer than HDT to build. 
These are measured format and implementation choices, not an intrinsic requirement of semantic conversion.

The v3.1.0 whole-genome benchmark took 16.07\,h including both optional builds: 20,911\,s for HDT and 31,230\,s for COTTAS, approximately 14.5\,h together. 
Retaining HDT, its index, and COTTAS used a reported 10.88\,GB; gzip-framed N-Triples alone occupied 2.89\,GB. 
These are alternative artifact choices, not a single minimum storage footprint. 
They must not be conflated with the later, differently configured 54-minute use-case conversion (with ClinVar). 
Supplementary Section~\ref{S-sup:measurements} retains per-input sizes, build times, and workspace measurements.

\subsubsection{Indexed retrieval benefits some questions, not every workflow}
For individual questions on 17.1M triples, indexed QLever took 0.005--4.62\,s against cyvcf2 scans of 11.5--12.3\,s (Figure~\ref{fig:retrieval}). 
Yet QLever first needed a 22.8-second index, in addition to conversion. 
For the complete thirteen-question batch, one parser pass took 12.4\,s, compared with 17.1\,s of QLever queries before setup. 
Thus the tested single cold questions and one-pass summaries favored the parser; selected repeated questions favored the existing graph index.

\begin{figure}[!htbp]
\centering
\includegraphics[width=\textwidth]{fig-retrieval.pdf}
\caption{
  \textbf{Retrieval cost depends on the question and execution path.} 
  Answers agree before timings are compared. 
  (a) Individual questions and the batch on 17.1M triples, mean of three replicates; conversion is excluded and QLever indexing appears only on the batch row. 
  (b) Four engines on a different, 0.96M-triple graph, excluding setup; means and ranges over nine runs. 
  (c) QLever on 17.1M triples loaded from three artifacts. 
  Similar times in (c) describe QLever's index, not native querying of those formats.
  }\label{fig:retrieval}
\end{figure}

Regional retrieval provides the stronger indexed-VCF baseline (Table~\ref{tab:regional}). 
Across both graphs, 11,160 executions had no failures or answer disagreements. 
Indexed VCF was faster at 1\,kb; QLever was faster for the tested 10\,Mb windows. 
Its observed cost was nearly flat across these windows, whereas indexed-reader costs increased with the records processed. 
Setup remained unequal: 23.2\,s for QLever versus 0.48\,s for bgzip and tabix.

\begin{table}[!htbp]
\caption{
  Regional retrieval on 17.1M triples, excluding setup. 
  Entries are medians in milliseconds over five questions' window medians and three repetitions. 
  Indexed arms use twenty windows per size; the unindexed scan uses three. 
  All arms use the same POS-based selection.
  }\label{tab:regional}
\centering\small
\setlength{\tabcolsep}{7pt}\renewcommand{\arraystretch}{1.13}
\begin{tabular}{@{}l r r r r@{}}
\toprule
\textbf{Execution path} & \textbf{1 kb} & \textbf{100 kb} & \textbf{1 Mb} & \textbf{10 Mb}\\
\midrule
QLever & 9.6 & 10.0 & 9.6 & 11.1\\
cyvcf2, indexed & 3.3 & 4.8 & 14.9 & 117.8\\
bcftools, indexed & 4.3 & 4.8 & 7.5 & 32.7\\
cyvcf2, unindexed & 635.2 & 638.0 & 632.0 & 649.9\\
\bottomrule
\end{tabular}
\end{table}

At whole-genome scale, QLever answered all thirteen inspection questions in 634.6\,s after an 868.0-second index build, with every answer equal to the parser's. 
These are separate retrieval checks without SHACL. 
Loading the same graph from HDT or COTTAS changed subsequent query time by only 0.3\%, but decoding costs differed. 
The native HDT- and COTTAS-backed paths did not complete at 171M triples, despite answering five questions. 
Supplementary Table~\ref{S-tab:scale} reports both the completions and failures; small-fixture engine rankings do not establish large-graph feasibility.

\section{Discussion}\label{sec4}
\subsection{The benefit is reusable meaning, not a universally better parser}\label{subsec:when}
The successful bcftools comparator is important: the integration task is not exclusive to RDF. 
What VCF-RDFizer contributes is a common place to state and inspect the relationships on which the task depends. 
Source observations retain their own context, external classifications and frequencies can be managed separately, and release rules can target declared identities or selections. 
The frequency join and the identity-targeted policy make this distinction concrete. 
They demonstrate a reusable implementation of the tested workflow, not that conventional pipelines cannot express it.

The policy plug-in adds a capability that file-level access control lacks: from one converted cohort, each requester receives a view of the genomic data materialized from the policies attached to it, and each view is checked against those policies and the VCF text before use. 
Because policies can target files, declared selections, or shared variant identities, the same mechanism can withhold a withdrawn participant, a gene panel, or a single variant without editing the source files. 
In this study the views differed by participant and gene panel, a single-variant rule was tested on a change fixture, and the whole-genome views were all-or-nothing.

Shared identifiers nevertheless require a normalization contract. 
Minimally trimmed SPDI expressions are not necessarily NCBI's contextual or canonical alleles in repeat regions \cite{holmes2020spdi}. 
An equivalent allele with a different identifier can miss an equality join and an identity-targeted rule. 
The demonstration establishes contig-alias handling and agreement after the stated preprocessing; broader identity equivalence remains a boundary to test. 
Likewise, simulated consent establishes executable rule behavior, not the validity of real consent or protection against access outside the evaluator. 
Multi-sample masking, enforcement of obligations, and secure deployment remain separate work.

\subsection{Validation needs to test both the graph and its oracle}\label{subsec4.1}
The strongest validation result is not simply a high pass count. 
It is the distinction between what source comparisons see, what structural constraints see, and what neither current implementation sees at scale. 
The injected faults expose that division: neither the query comparisons nor the default shapes detected all of them, and only the costly deeper profiles closed the gap. 
Their coverage therefore cannot yet be assumed for routine large-file validation, and the default profile's clean result on 58 ordinary validations is a conformance result, not a fault-detection rate. 
The real defects show why dependencies matter: the checks found three during development, none of them visible to a record count, which is the practical case for shipping validation with the converter. 
The consumer-genome follow-up adds another lesson: a valid conversion can fail when the oracle omits a modeled structure or relies on a different lexical form for an equal value.

The immediate priorities are therefore explicit. 
Extend census expectations to phase sets and other emitted structures, make value digests insensitive to permitted numeric serialization differences, and implement affordable checks for the dependencies now covered only by deep shapes. 
Until then, the 984 applicable campaign passes, the separate whole-genome summary agreement, and the 58.2M-triple diagnostic comparison must remain distinct claims. 
No experiment establishes exhaustive semantic equivalence or the biological correctness of the original calls.

\subsection{Representation and workload determine the deployment choice}\label{subsec4.2}
Expanded and condensed output exchange directly addressable structure for compact vectors. 
That affects queries, validation, and policy targeting as well as storage. 
Condensed sample-level checking needs vector-aware logic, and withholding an individual value would require rewriting the vector; the implemented release demonstration does neither. 
Triple counts alone are consequently an inadequate measure of the cost or capability of a genotype representation.

For a one-pass analysis or a narrow indexed lookup, the measured native VCF paths avoid setup that the RDF route cannot recover in that task. For repeated, linked, or policy-aware questions, gzip-framed N-Triples and an appropriate graph index are sufficient; optional HDT/COTTAS archives should be built only when their storage or access properties are needed. 
The 44-question break-even point is one scenario, not a tool constant: it includes a 423.7-second conversion workflow with HDT plus 22.8\,s of indexing, divided by about 10.2\,s saved per separate question. 
It assumes the measured question mix and repeated parser scans. 
A batched parser pass changes the comparison, and the minimal RDF-only setup was not isolated by that estimate (Supplementary Section~\ref{S-sup:measurements}).

\subsection{Scope and next steps}\label{subsec4.3}
The evidence supports a software and workflow contribution, not clinical deployment, measured usability, or superiority over other RDF converters. 
The converter comparison is documentation-based; the task comparison uses bcftools. 
Sample scaling holds INFO fixed, the whole cohort was not converted, and full-shape validation was limited to small inputs. 
The use case has single-sample files and simulated policies. 
Release-specific results must also remain tied to the versions that produced them.

Focused next steps are more useful than another broad benchmark: resolve the observed oracle failures, complete a conformant sites-only validation, test equivalent variant representations against identity-targeted policies, and assess selective non-empty release at whole-genome scale. 
A shared-input comparison with another converter would strengthen comparative claims. 
GA4GH VRS identifiers, federated queries, and multi-sample policy evaluation are extensions, not demonstrated capabilities of this study.

\Needspace{9\baselineskip}
\section{Conclusions}\label{sec:conclusions}
VCF-RDFizer makes VCF information explicit in a representation that can be checked against its source, connected to external knowledge, and used in policy-aware research queries. 
The evaluation demonstrates complementary validation, reusable joins, and requester-specific result agreement with a conventional workflow, while exposing limits in oracle coverage, identity assumptions, and large-graph validation. 
Its measured costs support choosing semantic conversion for tasks that need those relationships and checks, rather than treating RDF as a universal replacement for native VCF processing.

\section*{Availability and requirements}
\textbf{Project:} VCF-RDFizer. 
\textbf{Home page:} \url{https://github.com/ecrum19/VCF-RDFizer}. 
\textbf{Language:} Python, with the mapping, compression, and query tools specified in Supplementary Section~\ref{S-sup:repro}. 
\textbf{Platforms:} Linux, macOS, and Windows are covered by the tool's tests; the experiments reported here ran on Ubuntu. 
\textbf{Distribution:} PyPI, Conda-Forge, and Docker. 
\textbf{License:} MIT. \textbf{Any restrictions to use by non-academics:} none. Docker supplies the pinned benchmark environment; requirements for other installations are documented with each release. 
v3.1.0 produced the base benchmark, v3.2.0 converted the use-case inputs, and v3.3.0 releases the linking and policy extensions evaluated in pre-release builds. 
The release ledger distinguishes those builds rather than attributing every result to one version.

\section*{Declarations}
\subsection*{Ethics approval and consent to participate}
This study used publicly released genomic data and involved no recruitment of or interaction with participants. 
Personal Genome Project donors released their files under the project's open-consent model \cite{personal_genome_proj}; GIAB and precisionFDA files are public reference materials \cite{nist_gib_2025,giab_v421_2022}; the 1000 Genomes and HGSVC2 datasets are released for research reuse \cite{thousand_genomes_2015,byrska2022highcoverage,hgsvc_2021}; and ClinVar is a public NCBI resource \cite{clinvar_vcf}. 
The experimental consents are simulated and do not express any real participant's wishes. 
No identification of participants was attempted. Ethics approval was not sought for the reported work.

\subsection*{Consent for publication}
Not applicable. 
The manuscript reports aggregate results and software measurements, not individual-level genomic records, images, or personal details.

\subsection*{Availability of data and materials}
Input filenames, provider URLs and download commands, fixtures, derivation scripts, query files, manifests, normalized outputs, and comparison reports are maintained at \url{https://github.com/ecrum19/vcf-rdfizer-testing}. 
Large source VCFs are obtained from their public providers rather than redistributed in Git. 
Experiment~17 (\texttt{benchmarks/use\_case/acmg}) contains the use-case definitions, simulated policies, conventional baseline, and applied-change tests. 
Its match sets, timings, and check reports, as well as the follow-up validation runs, are recorded under \texttt{BioMedSem\_2026/benchmark-results}. 
No biological materials were generated or used.

The software is available at the project home page. 
The base benchmark used v3.1.0 (\texttt{d3b34d5}); the linking and policy extensions are released in v3.3.0 (\texttt{7ab2300}; image \texttt{sha256:1838720d\ldots06432}); the use case ran pre-release builds (\texttt{578c2be} for the whole genome), which v3.3.0 changes only by refusing an empty selector list and an RDF input without VCF triples. 
Supplementary Section~\ref{S-sup:repro} identifies the pre-release and later-run commits, the v3.1.0 image digest, and the bundled shape snapshot. 
The public vocabulary is separately versioned \cite{vcfc210}. 
Results must be traced to their recorded build and input, not only to the latest package.
\authorquery{Add a persistent archive identifier containing the exact evaluated code, dirty-tree changes for use-case arms 1--2, complete image digests, vocabulary/shapes, queries, and evidence manifests. Record the precise commits and configurations for the new default-shape mutation and NG131FQA1I follow-up runs. The supplied manuscript package does not itself establish this complete archived snapshot.}

\subsection*{Competing interests}
None declared.

\subsection*{Funding}
Elias Crum acknowledges support from the Research Foundation--Flanders (FWO), SB Fellowship 1S27825N.
\authorquery{Complete the funding statements for B.B., R.T., and G.E.}

\subsection*{Authors' contributions}
E.C., G.E., and R.T. conceived of this study. 
E.C. did the implementation and coding work. 
E.C. devised the validation and use-case demonstration with the guidance of R.T. and executed the experiments. 
E.C., G.E., and R.T. reviewed and verified all results. 
E.C., B.B., G.E., and R.T. wrote and edited the manuscript.

\subsection*{Acknowledgments}
Benchmarking used the resources of IDLab at Ghent University. 
Claude Opus 5.5~\cite{anthropic2026claudeopus55} and GPT-6 Astra~\cite{openai2026gpt6astra} assisted with experiment automation scripting, result data analysis scripting, and restructuring and language editing of this manuscript. The authors reviewed all AI outputs and take full responsibility for the manuscript, experiments, tools, and repositories.

\clearpage
\bibliography{reference}
\end{document}
